\documentclass[a4paper,12pt]{article}
\usepackage[utf8]{inputenc}
\usepackage[T1]{fontenc}
\usepackage[ngerman]{babel}
\usepackage{amsmath}
\usepackage{amssymb}
\usepackage{enumitem}
\usepackage{array}
\usepackage{geometry}
\geometry{a4paper, margin=2.5cm}
\usepackage{fancyhdr}
\usepackage{parskip}

% Punkteangabe rechtsbündig am Ende eines Aufgabenteils
\newcommand{\pkt}[1]{\hfill\textit{(#1\,P)}}
% Kopf einer Aufgabe: \aufgabe{Nummer}{Titel}{Gesamtpunkte}
\newcommand{\aufgabe}[3]{%
  \bigskip\par\noindent\textbf{Aufgabe #1: #2}\hfill\textbf{#3 Punkte}%
  \par\nobreak\medskip}

\pagestyle{fancy}
\fancyhf{}
\lhead{\small 61211 Analysis -- Probeklausur 3}
\rhead{\small Matr.-Nr.: \makebox[3cm]{\dotfill}}
\cfoot{\small Seite \thepage}
\renewcommand{\headrulewidth}{0.4pt}

\begin{document}
\thispagestyle{empty}

\begin{center}
  {\LARGE\bfseries Probeklausur 3}\\[0.4em]
  {\large 61211 Analysis}\\[0.6em]
  {\normalsize FernUniversität in Hagen $\cdot$ Fakultät für Mathematik und Informatik}
\end{center}

\vspace{0.8em}
\noindent\small\textit{Diese Probeklausur ist bewusst beweis- und begriffsorientiert
gehalten. Statt reiner Rechenaufgaben stehen \glqq beweisen\grqq{} und \glqq beweisen
oder widerlegen\grqq{} im Vordergrund (Ordnungsvollständigkeit, Zwischenwertsatz,
Mittelwertsatz, Fixpunktsätze, Taylor, Approximation, Kurvenintegrale). Alle Aufgaben
lassen sich mit den Sätzen und Definitionen des Kurstextes lösen.}\normalsize

\vspace{1.0em}

\noindent
\begin{tabular}{@{}ll@{}}
  Name: & \makebox[7cm]{\dotfill}\\[1.0em]
  Vorname: & \makebox[7cm]{\dotfill}\\[1.0em]
  Matrikelnummer: & \makebox[7cm]{\dotfill}\\
\end{tabular}

\vspace{1.2em}

\begin{center}
  \renewcommand{\arraystretch}{1.5}
  \begin{tabular}{|l|c|c|c|c|c|c|c||c|}
    \hline
    \textbf{Aufgabe} & 1 & 2 & 3 & 4 & 5 & 6 & 7 & \textbf{$\Sigma$}\\
    \hline
    max.\ Punkte & 20 & 12 & 14 & 16 & 14 & 14 & 10 & 100\\
    \hline
    erreicht & & & & & & & & \\
    \hline
  \end{tabular}
\end{center}

\vspace{0.8em}

\noindent\textbf{Bearbeitungshinweise}
\begin{itemize}[leftmargin=1.5em, itemsep=2pt]
  \item Die Bearbeitungszeit beträgt \textbf{120 Minuten}. Es sind alle 7 Aufgaben zu bearbeiten.
  \item Zugelassenes Hilfsmittel: ein handbeschriebenes DIN-A4-Blatt (auch beidseitig) mit eigenen Notizen.
  \item Verwenden Sie keinen Bleistift und keinen Rotstift.
  \item Formulieren Sie Ihre Beweise in vollständigen Sätzen und begründen Sie jeden Schritt. Verweise auf Sätze des Kurstextes dürfen benutzt werden.
\end{itemize}

\newpage

% =====================================================================
\aufgabe{1}{Multiple Choice}{20}

Kreuzen Sie an, ob die jeweilige Aussage \emph{wahr} oder \emph{falsch} ist. Für jede
richtig beantwortete Aussage erhalten Sie $2$ Punkte, für jede unbeantwortete Aussage
$1$ Punkt und für jede falsch beantwortete Aussage $0$ Punkte.

\medskip
\renewcommand{\arraystretch}{1.5}
\noindent\begin{tabular}{@{}c@{\hspace{0.6em}}p{11.4cm}@{\hspace{0.8em}}c@{\hspace{1.0em}}c@{}}
  & & \small Wahr & \small Falsch\\[0.2em]
  (a) & Jede nichtleere, nach oben beschränkte Teilmenge von $\mathbb{Q}$ besitzt ein Supremum in $\mathbb{Q}$. & $\square$ & $\square$\\
  (b) & Der normierte Raum $(C([0,1]),\|\cdot\|_1)$ mit $\|f\|_1=\int_0^1|f(t)|\,dt$ ist vollständig. & $\square$ & $\square$\\
  (c) & Die Menge $\mathbb{Q}\cap[0,1]$ ist kompakt in $\mathbb{R}$. & $\square$ & $\square$\\
  (d) & Ist $M\subseteq\mathbb{R}^n$ kompakt und $f:M\to\mathbb{R}^m$ stetig, so ist $f(M)$ kompakt. & $\square$ & $\square$\\
  (e) & Ist $f:\mathbb{R}^n\to\mathbb{R}$ in $a$ differenzierbar, so existiert für jeden Einheitsvektor $v$ die Richtungsableitung $D_v f(a)$, und es gilt $D_v f(a)=\operatorname{grad} f(a)\cdot v$. & $\square$ & $\square$\\
  (f) & Ist $f:\mathbb{R}^n\to\mathbb{R}^n$ stetig differenzierbar mit $\operatorname{Rang} f'(a)=n$, so ist $f$ auf ganz $\mathbb{R}^n$ injektiv. & $\square$ & $\square$\\
  (g) & Ist $f$ zweimal stetig differenzierbar, $\operatorname{grad} f(a)=0$ und die Hesse-Matrix $H(a)$ positiv definit, so besitzt $f$ in $a$ ein lokales Minimum. & $\square$ & $\square$\\
  (h) & Jede monotone Funktion $f:[a,b]\to\mathbb{R}$ (mit $a<b$) ist Riemann-integrierbar. & $\square$ & $\square$\\
  (i) & Für die Fourierkoeffizienten $(a_k),(b_k)$ einer Funktion $f\in\mathcal{R}(\,]-\pi,\pi])$ gilt $\lim_{k\to\infty}a_k=\lim_{k\to\infty}b_k=0$. & $\square$ & $\square$\\
  (j) & Jede stetige Kurve $\varphi:[a,b]\to\mathbb{R}^n$ ist rektifizierbar. & $\square$ & $\square$\\
\end{tabular}

% =====================================================================
\aufgabe{2}{Supremum und Infimum}{12}

Seien $A,B\subseteq\mathbb{R}$ nichtleer und beschränkt.
\begin{enumerate}[label=\alph*)]
  \item Beweisen Sie $\displaystyle\sup(A\cup B)=\max\{\sup A,\ \sup B\}$. \pkt{4}
  \item Für $A+B:=\{a+b\mid a\in A,\ b\in B\}$ beweisen Sie
  \[ \sup(A+B)=\sup A+\sup B. \]
  \pkt{5}
  \item Für $A\cdot B:=\{a\,b\mid a\in A,\ b\in B\}$: Beweisen oder widerlegen Sie, dass \emph{stets}
  $\sup(A\cdot B)=\sup A\cdot\sup B$ gilt. \pkt{3}
\end{enumerate}

% =====================================================================
\aufgabe{3}{Zwischenwertsatz und Fixpunkte}{14}

\begin{enumerate}[label=\alph*)]
  \item Sei $f:[0,1]\to[0,1]$ stetig. Beweisen Sie, dass $f$ einen \emph{Fixpunkt} besitzt,
  d.\,h.\ es gibt ein $\xi\in[0,1]$ mit $f(\xi)=\xi$. \pkt{6}
  \item Beweisen oder widerlegen Sie: Jede stetige Funktion $f:\,]0,1[\,\to\,]0,1[$ besitzt einen Fixpunkt. \pkt{4}
  \item Beweisen Sie mit Hilfe des Zwischenwertsatzes, dass jedes Polynom \emph{ungeraden} Grades
  mit reellen Koeffizienten mindestens eine reelle Nullstelle besitzt. \pkt{4}
\end{enumerate}

% =====================================================================
\aufgabe{4}{Mittelwertsatz und Kontraktionen}{16}

Sei $f:\mathbb{R}\to\mathbb{R}$ differenzierbar.
\begin{enumerate}[label=\alph*)]
  \item Es gelte $q:=\sup_{x\in\mathbb{R}}|f'(x)|<1$. Beweisen Sie, dass $f$ dehnungsbeschränkt
  mit Lipschitzkonstante $q$ ist, d.\,h.
  \[ |f(x)-f(y)|\le q\,|x-y|\qquad\text{für alle }x,y\in\mathbb{R}. \]
  \pkt{5}
  \item Folgern Sie mit dem Banachschen Fixpunktsatz (Satz~4.1.6), dass $f$ genau einen Fixpunkt
  $x^\ast\in\mathbb{R}$ besitzt und dass die Iterationsfolge $x_k:=f(x_{k-1})$ für jeden Startwert
  $x_0\in\mathbb{R}$ gegen $x^\ast$ konvergiert. \pkt{6}
  \item Beweisen oder widerlegen Sie: Ist $f:\mathbb{R}\to\mathbb{R}$ differenzierbar mit $|f'(x)|<1$
  für \emph{jedes} $x\in\mathbb{R}$, so besitzt $f$ (mindestens) einen Fixpunkt. \pkt{5}
\end{enumerate}

% =====================================================================
\aufgabe{5}{Konvexität, Taylor und Extrema}{14}

\begin{enumerate}[label=\alph*)]
  \item Sei $f:\mathbb{R}\to\mathbb{R}$ zweimal stetig differenzierbar mit $f''(x)\ge 0$ für alle
  $x\in\mathbb{R}$. Beweisen Sie mit dem Satz von Taylor (Satz~4.3.9), dass für alle $a,x\in\mathbb{R}$
  \[ f(x)\ge f(a)+f'(a)\,(x-a) \]
  gilt (der Graph liegt oberhalb jeder Tangente). \pkt{6}
  \item Folgern Sie: Gilt zusätzlich $f'(a)=0$, so besitzt $f$ in $a$ ein \emph{globales} Minimum. \pkt{3}
  \item Beweisen oder widerlegen Sie: Ist $f:\mathbb{R}^2\to\mathbb{R}$ zweimal stetig differenzierbar
  mit $\operatorname{grad} f(a)=0$ und ist die Hesse-Matrix $H(a)$ positiv \emph{semidefinit}, so besitzt
  $f$ in $a$ ein lokales Minimum. \pkt{5}
\end{enumerate}

% =====================================================================
\aufgabe{6}{Stetige Funktionen und Approximation}{14}

Sei $I=[a,b]$ mit $a<b$.
\begin{enumerate}[label=\alph*)]
  \item Sei $f:I\to\mathbb{R}$ stetig mit $\displaystyle\int_I f(x)^2\,dx=0$. Beweisen Sie, dass
  $f\equiv 0$ ist (d.\,h.\ $f(x)=0$ für jedes $x\in I$). \pkt{6}
  \item Sei $f:I\to\mathbb{R}$ stetig, und es gelte $\displaystyle\int_I f(x)\,p(x)\,dx=0$ für
  \emph{jede} Polynomfunktion $p$. Beweisen Sie $f\equiv 0$.
  \emph{Hinweis:} Weierstraßscher Approximationssatz (Satz~5.5.4); verwenden Sie a). \pkt{8}
\end{enumerate}

% =====================================================================
\aufgabe{7}{Kurvenlänge und Wegunabhängigkeit}{10}

\begin{enumerate}[label=\alph*)]
  \item Sei $W$ eine glatte Kurve in $\mathbb{R}^n$ mit glatter Parameterdarstellung
  $\varphi:[\alpha,\beta]\to\mathbb{R}^n$, Anfangspunkt $p:=\varphi(\alpha)$ und Endpunkt
  $q:=\varphi(\beta)$. Beweisen Sie
  \[ L(W)\ \ge\ \|q-p\|_2 \]
  (die geradlinige Strecke ist die kürzeste Verbindung).
  \emph{Hinweis:} $q-p=\int_\alpha^\beta \dot\varphi(t)\,dt$ (komponentenweise) und
  Cauchy-Schwarzsche Ungleichung. \pkt{5}
  \item Betrachtet werde das Vektorfeld
  \[ f:\mathbb{R}^2\setminus\{0\}\to\mathbb{R}^2,\qquad
     f(x,y):=\Big(\tfrac{-y}{x^2+y^2},\ \tfrac{x}{x^2+y^2}\Big). \]
  Beweisen oder widerlegen Sie: $f$ erfüllt zwar die Integrabilitätsbedingungen
  $D_1 f_2=D_2 f_1$, besitzt aber \emph{keine} Stammfunktion auf $\mathbb{R}^2\setminus\{0\}$. \pkt{5}
\end{enumerate}

\vfill
\begin{center}\small Viel Erfolg!\end{center}

\end{document}
